
\documentclass[12pt,onecolumn]{article}

\usepackage[utf8]{inputenc}
\usepackage[T1]{fontenc}
\usepackage{lmodern}
\usepackage[margin=1in]{geometry}
\usepackage{setspace}
\usepackage{microtype}
\usepackage{amsmath,amssymb,amsthm}
\usepackage[numbers,sort&compress]{natbib}
\usepackage[colorlinks=true,linkcolor=blue,citecolor=blue,urlcolor=blue]{hyperref}
\usepackage{orcidlink}
\usepackage{fancyhdr}

\setstretch{1.08}
\setlength{\headheight}{14pt}

\pagestyle{fancy}
\fancyhf{}
\fancyhead[L]{Space Is Not a Location}
\fancyhead[R]{John C. W. McKinley}
\fancyfoot[C]{\thepage}

\newtheorem{proposition}{Proposition}
\newtheorem{definition}{Definition}

\title{\textbf{Space Is Not a Location.\\
It Is a Relationship.}\\[0.5em]
\large A Relational Ontology of Distance Without
an Enclosing Universe}

\author{
John C. W. McKinley\\
Independent Researcher\\
\orcidlink{0009-0005-7097-5035}\\
\href{https://orcid.org/0009-0005-7097-5035}
{ORCID: 0009-0005-7097-5035}
}

\date{October 2026}

\begin{document}

\maketitle

\begin{abstract}
Space is conventionally represented as a geometric
arena within which physical objects and events are
located. This paper challenges the ontological
assumption that space exists independently as a
container for the universe. A container requires
boundaries, an interior, and an exterior. If space
itself is treated as a container, the question
immediately arises: where is that container located?
Placing it within a larger space merely reproduces
the original problem, generating an infinite regress
of enclosing spatial structures. Rejecting that
regress motivates a simpler interpretation: space
is not an independently existing location but a
system of measurable relationships among physical
entities and events.

Under this interpretation, distance does not describe
separation within a preexisting spatial substance.
Rather, distance is the quantitative expression of
spatial relationships. Geometry describes the
organization of those relationships, not the shape
of an invisible container. The universe therefore
requires no external location, surrounding
environment, or enclosing structure. The question
of where the universe exists becomes a category
error: location is a relationship defined within
the universe, not a property that the universe must
possess relative to something outside itself.

This paper develops the argument as an ontological
proposition compatible with the measurable geometry
of special and general relativity. It distinguishes
the empirical reality of spatial measurements from
the presumed independent existence of space as a
physical entity. The resulting thesis is direct:
\emph{space is not something that exists somewhere.
Space is the relational structure by which somewhere
is defined.}
\end{abstract}

\noindent\textbf{Keywords:}
relational space, ontology of space, spatial geometry,
infinite regress, general relativity, relational
distance, spacetime foundations

\section{Introduction}

Where is the universe?

The question appears straightforward. Every familiar
physical object occupies a location. A house occupies
land. A planet occupies a region of the solar system.
A galaxy occupies a position relative to other
galaxies. It therefore seems natural to ask where
the universe itself is located.

But the question contains an assumption that does
not follow from the preceding examples.

A house is located relative to its surroundings.
A planet is located relative to other celestial
bodies. A galaxy is located relative to a larger
distribution of matter.

In each example, location is relational.

The universe, understood as the totality of physical
reality, has no stipulated external physical
reference system. Asking where it is located
therefore demands a relationship that has not
been defined.

The difficulty arises when the familiar experience
of objects occupying space is extended into the
assumption that space itself must occupy something.

This paper rejects that extension.

Its central proposition is:

\begin{quote}
\emph{Space is not a location. Space is the system
of relationships through which location and
distance are defined.}
\end{quote}

The argument is ontological rather than a proposed
modification of established physical measurements.

Distances remain measurable. Geometry remains
applicable. Relativistic intervals remain physically
significant.

What is rejected is the additional assumption that
these relationships require an independently existing
spatial container.

\section{The Container Assumption}

\subsection{What Is a Container?}

A physical container is an object or structure
distinguished from its surroundings.

It possesses an interior, a boundary, and an exterior,
at least in the ordinary sense in which physical
containers are understood.

A room has walls. A bottle has a surface. A box
occupies a region within a larger environment.

These examples share a common feature: the
container is itself spatially situated.

The concept of a container therefore presupposes
spatial relationships.

It does not explain their origin.

\subsection{The Box Inside a Box}

Suppose the universe is imagined as occupying
a container $C_0$.

If $C_0$ is a physical container of the ordinary
kind, it has an exterior and is situated relative
to a surrounding environment.

Call that environment $C_1$.

If $C_1$ is also understood as a container,
the same question arises.

We obtain the sequence

\begin{equation}
C_0 \subset C_1 \subset C_2
\subset \cdots .
\end{equation}

The sequence does not explain spatial existence.
It repeatedly assumes the very thing it was
introduced to explain.

The issue is not that an infinite mathematical
sequence is intrinsically contradictory.

The issue is explanatory.

Each successive container inherits the spatial
assumptions of its predecessor.

No member of the sequence explains why spatial
relationships exist in the first place.

\subsection{The Holodeck Analogy}

A holodeck is an especially revealing analogy.

A simulated environment appears to contain
distances, objects, and locations.

But the simulation is produced by machinery
situated in a larger physical environment.

The holodeck therefore possesses two distinct
levels:

\begin{enumerate}
\item The spatial relationships represented
inside the simulation.
\item The external physical environment containing
the machinery.
\end{enumerate}

To describe the universe as a holodeck is to
introduce an external environment by analogy.

That environment is not established by the
observed geometry of the universe.

It is an additional assumption.

The universe need not be a simulation housed
inside a larger room merely because familiar
simulations require hardware.

\section{The Regress Argument}

The container problem can be stated formally.

\begin{proposition}
A theory that explains every spatial container
solely by placing it within another spatial
container does not provide a terminating
explanation of spatial containment.
\end{proposition}

\begin{proof}
Let $C_n$ denote a spatial container.

Assume the proposed explanation of $C_n$
requires a distinct enclosing spatial container
$C_{n+1}$.

Then explaining $C_0$ requires $C_1$.

Explaining $C_1$ requires $C_2$.

By repeated application, explaining any
$C_n$ requires $C_{n+1}$.

Consequently, the explanatory rule supplies
no terminal member.

If a terminating explanation is required,
the rule must eventually be abandoned or
replaced.

Thus, nesting alone does not furnish a
terminating explanation of spatial containment.
\end{proof}

This proposition is conditional.

It does not prove that every infinite regress
is mathematically impossible, nor does it
exclude every conceivable fundamental spatial
structure.

It establishes that a chain of containers
cannot explain spatial containment merely
by repeating containment.

The proposed solution is to reject the initial
ontological assumption.

Space need not be a container at all.

\section{Space as Relationship}

\subsection{The Primary Definition}

Consider two physical entities, $A$ and $B$.

Their spatial relationship is characterized
by a measurable separation.

In an appropriate spatial geometry, write

\begin{equation}
d(A,B).
\end{equation}

The distance is a relationship between
the entities.

It is not necessary to introduce a third
physical object called ``space'' occupying
the interval between them.

The distinction is between the mathematical
structure used to describe separation and
the assertion that this structure is itself
a substance.

\begin{definition}
In the relational interpretation adopted here,
space is the organized structure of spatial
relations among physical entities and events,
rather than an independently existing material
container.
\end{definition}

This definition preserves the physical
significance of distance.

It changes what is taken to exist independently.

\subsection{Distance Without Substance}

Suppose two objects are separated by one meter.

The measurement has operational meaning.

A measuring rod, a calibrated signal procedure,
or an equivalent geometric construction
establishes the separation.

But the measurement does not, by itself,
establish that a meter of independently
existing spatial substance occupies the
interval.

The relational interpretation asserts:

\begin{equation}
\boxed{
\text{Spatial distance}
\neq
\text{spatial substance}.
}
\end{equation}

Distance is a measurable property of a
physical configuration.

It is not automatically a separate physical
object.

\subsection{A Network of Relations}

For a collection of entities

\begin{equation}
\mathcal{P}=\{P_1,P_2,\ldots,P_N\},
\end{equation}

a relational description may be represented
by a set of separations

\begin{equation}
\mathcal{R}
=
\left\{
d(P_i,P_j)
\right\}_{i,j}.
\end{equation}

For an ordinary metric distance function,

\begin{align}
d(P_i,P_j)&\geq 0,\\
d(P_i,P_j)&=d(P_j,P_i),\\
d(P_i,P_k)&\leq
d(P_i,P_j)+d(P_j,P_k).
\end{align}

With the additional condition that distinct
points have positive separation, these are
the familiar metric-space properties.

They demonstrate how geometric organization
can be represented through relationships.

The mathematical points need not be imagined
as beads suspended in a physical box.

They are elements of a structure whose
relationships are specified.

This construction is illustrative rather
than a complete reconstruction of relativistic
spacetime.

In particular, the spatial distances of
relativity depend upon the choice of a
simultaneity convention or spacelike
hypersurface.

The invariant spacetime interval requires
a Lorentzian rather than an ordinary
positive-definite metric.

\section{Location Is Not an Intrinsic Property}

The relational interpretation distinguishes
an entity from its assigned coordinates.

Consider an object $A$ assigned coordinates

\begin{equation}
x^\mu_A.
\end{equation}

These coordinates describe the object relative
to a chosen coordinate system.

A change of coordinates,

\begin{equation}
x^\mu \longrightarrow x'^\mu(x),
\end{equation}

changes the numerical description without
necessarily changing the physical situation.

This is familiar in general relativity,
where coordinate labels do not by themselves
constitute observable physical quantities
\citep{einstein1916,wald1984}.

A coordinate is not a physical address
written into the universe independently
of all reference conventions.

A location becomes physically meaningful
through relationships, measurements,
and identifiable events.

Thus the question

\begin{quote}
Where is the universe?
\end{quote}

cannot be answered by assigning the universe
a coordinate within itself.

Nor does the question automatically justify
postulating a larger coordinate system
outside physical reality.

A coordinate system is a descriptive
construction, not evidence of an external
container.

\section{Compatibility with Special Relativity}

Special relativity provides an important
test of the relational interpretation.

For two events in Minkowski spacetime,
the invariant interval, using signature
$(-,+,+,+)$, is

\begin{equation}
ds^2
=
-c^2dt^2+dx^2+dy^2+dz^2.
\end{equation}

Different inertial observers assign different
coordinate times and spatial separations
to the same pair of events.

Lorentz transformations preserve the
spacetime interval.

For motion along the $x$ direction,

\begin{align}
x' &= \gamma(x-vt),\\
t' &= \gamma
\left(t-\frac{vx}{c^2}\right),
\end{align}

where

\begin{equation}
\gamma
=
\frac{1}{\sqrt{1-v^2/c^2}}.
\end{equation}

These transformations demonstrate that
spatial separation and temporal separation
are not universally fixed coordinate
quantities.

Their values depend on the observer's
inertial frame.

The relational interpretation takes this
dependence seriously.

It does not require a privileged material
container in which all observers must
agree upon absolute positions.

Instead, physical relationships are
expressed through frame-dependent
coordinates and invariant structures.

Special relativity does not, by itself,
prove relational ontology.

It does show that a physically successful
description does not require Newtonian
absolute space as an observable reference
frame.

\section{Compatibility with General Relativity}

General relativity replaces the fixed
Minkowski metric with a dynamical
Lorentzian metric,

\begin{equation}
ds^2
=
g_{\mu\nu}dx^\mu dx^\nu.
\end{equation}

The metric determines local intervals,
causal structure, and the geometry
relevant to physical measurements.

Its dynamics are governed by the
Einstein field equations,

\begin{equation}
G_{\mu\nu}
+
\Lambda g_{\mu\nu}
=
\frac{8\pi G}{c^4}T_{\mu\nu}.
\end{equation}

These equations do not require the
four-dimensional manifold to be embedded
within a higher-dimensional physical
container.

The geometry is intrinsic.

Curvature is determined through
relationships defined on the manifold
itself.

For example, the Riemann curvature
tensor,

\begin{equation}
R^\rho{}_{\sigma\mu\nu},
\end{equation}

is constructed from the metric connection
and its derivatives.

It does not require an observer standing
outside the universe to inspect the
bending of an enclosing surface.

\subsection{Intrinsic Versus Extrinsic Geometry}

A two-dimensional surface can possess
intrinsic curvature.

The inhabitants of that surface could,
in principle, measure its geometry
without observing it from a third
spatial dimension.

Similarly, the curvature of relativistic
spacetime is not dependent upon a literal
external room.

The embedding of a geometric structure
into a higher-dimensional space is a
mathematical option, not an automatic
physical requirement.

This distinction is central to the
present argument.

A curved universe does not require
something outside it into which
it curves.

\subsection{The Ontological Distinction}

General relativity is compatible with
multiple philosophical interpretations
of spacetime.

Some interpretations treat spacetime
geometry as possessing independent
physical reality.

Others emphasize relational or
structural descriptions
\citep{rovelli2004,earman1989}.

The present paper does not claim that
Einstein's equations uniquely establish
relationalism.

It argues that their successful use
does not require space to be a
material container.

That distinction leaves the empirical
content of general relativity intact.

\section{The Universe Does Not Need an Outside}

The phrase ``outside the universe''
has no automatic physical definition.

If the universe denotes the totality
of physical reality, then a purported
physical exterior would itself belong
to that totality.

The concept of an exterior is meaningful
for bounded objects within a larger
spatial arrangement.

It does not follow that the same concept
applies to the totality of physical
relationships.

This can be expressed as a distinction
between two questions:

\begin{enumerate}
\item What are the spatial relationships
among physical entities?
\item What larger physical container
holds the totality of those relationships?
\end{enumerate}

The first question is meaningful
within physical geometry.

The second introduces an additional
structure not demanded by the first.

Under the proposed ontology,
the second question is unnecessary.

The universe is not an object
requiring an address.

It is the physical totality within
which addresses are defined.

\section{Material Boundaries Cannot Explain Space}

A material container is composed of
physical constituents.

Those constituents possess spatial
relationships among themselves.

The boundary of a box, for example,
is defined through the arrangement
of its material components.

Consequently, using a material box
to explain the existence of spatial
relationships reverses the explanatory
order.

The spatial relationships are already
presupposed in the construction of
the box.

\begin{proposition}
A material container cannot serve as
a non-circular explanation of spatial
relationships if its own definition
presupposes those relationships.
\end{proposition}

\begin{proof}
Let $C$ be a material container.

Its physical boundary and internal
structure are specified through
relationships among its constituents.

Therefore, the description of $C$
already assumes spatial relations.

An explanation that derives spatial
relations solely from the existence
of $C$ consequently assumes the
target of the explanation.

Such an explanation is circular.
\end{proof}

This proposition does not prohibit
a fundamental geometric structure
that is not a material box.

It rejects the specific claim
that an ordinary material container
explains the existence of spatial
relationships.

The relational alternative avoids
that circularity by treating the
relationships themselves as primary.

\section{Expansion Without an Expanding Box}

Cosmology provides a particularly
useful application.

In a homogeneous and isotropic
FLRW description, the metric is

\begin{equation}
ds^2
=
-c^2dt^2
+
a^2(t)
\left[
\frac{dr^2}{1-kr^2}
+
r^2d\Omega^2
\right],
\end{equation}

where $a(t)$ is the scale factor
and $k$ characterizes spatial
curvature.

For comoving observers at fixed
comoving separation $\chi$, their
proper separation on a constant
cosmic-time hypersurface is

\begin{equation}
D(t)=a(t)\chi.
\end{equation}

Differentiating gives

\begin{equation}
\dot D(t)=H(t)D(t),
\end{equation}

where

\begin{equation}
H(t)=\frac{\dot a(t)}{a(t)}.
\end{equation}

The physical content concerns the
evolution of measurable separations
among comoving observers.

It does not require the universe
to expand into a preexisting
external volume.

In the relational interpretation,
cosmic expansion describes the
changing metric relationships
among cosmological worldlines.

The question

\begin{quote}
What is the universe expanding into?
\end{quote}

therefore assumes an external
container that the equations
do not require.

The expansion is a change in
intrinsic geometry.

It is not the motion of an
outer wall through a larger room.

\section{A Universe Without a Spatial Edge}

An important distinction must be
made between finiteness and
the existence of a boundary.

A geometric space can have
finite volume without possessing
an edge.

For example, the three-sphere
$S^3$ is compact and has finite
volume, yet it has no boundary.

In standard geometry,

\begin{equation}
\partial S^3=\varnothing.
\end{equation}

This mathematical possibility
demonstrates that a finite
universe need not be imagined
as a finite box with an outer wall.

Likewise, a spatially infinite
cosmological model does not
require an infinite material
container.

Finite and infinite cosmological
geometries are both described
intrinsically.

The present thesis is independent
of whether the spatial universe
is finite or infinite.

Neither possibility requires
an enclosing physical room.

\section{Connection to the Timeless Light Model}

The relational interpretation
developed here is consistent
with the broader ontological
program of the Timeless Light
Model (TLM).

TLM distinguishes the Quantum
Platform (QP), understood as a
timeless instruction layer,
from the Spacetime Deployment
Frame (SDF), in which measurable
physical events and their
temporal relationships are
represented.

The present paper does not
depend upon the validity of
that framework.

Its central argument is
independent:

\begin{quote}
Measurable spatial relationships
do not establish the existence
of an independent spatial
container.
\end{quote}

Nevertheless, the distinction
is useful within TLM.

If spatial relationships
describe physical deployment,
the SDF need not be interpreted
as a literal physical box.

Its geometry characterizes
relationships among deployed
events, observers, and
physical systems.

This preserves the operational
reality of distance without
requiring space to be an
independent substance.

The QP should not, in turn,
be pictured as another room
outside the SDF.

Doing so would reproduce
the very container assumption
rejected in this paper.

The QP/SDF distinction is
ontological and explanatory,
not a claim that one spatial
box sits inside another.

\section{What the Argument Does and Does Not Establish}

The argument establishes
three principal conclusions.

First, the material-container
analogy cannot explain the
origin of spatial relationships
without presupposing them.

Second, placing the universe
inside progressively larger
spatial containers does not
resolve the explanatory problem.

Third, relational geometry
provides a coherent alternative
in which measurable distance
does not require an independently
existing material medium
called space.

The argument does not derive
the Einstein field equations,
the value of the speed of light,
or the observed cosmological
parameters.

It does not establish that
spacetime has no physical
degrees of freedom.

Gravitational waves, curvature,
and dynamical metric phenomena
remain part of established
general relativity.

Nor does the argument claim
that relationalism is a newly
discovered philosophical position.

Relational interpretations
of space have a long history,
including the Leibniz--Clarke
debate and subsequent work
in the foundations of physics
\citep{alexander1956,barbour1982}.

The contribution of this paper
is a compact container-regress
argument and its application
to a specific ontological thesis:
space is not a location but
a structure of relationships.

\section{Objections and Clarifications}

\subsection{Does Space Exist If Nothing Is Present?}

A relational ontology treats
spatial structure as grounded
in physical relations rather
than in an independently
existing container.

The question of an entirely
empty universe exposes a
difference between relational
and substantival interpretations.

A mathematical manifold
with a metric can be defined
without ordinary matter.

General relativity also
admits vacuum solutions.

The existence of such
solutions does not settle
whether the manifold represents
an independent substance.

The present thesis concerns
ontology, not the mathematical
availability of vacuum geometry.

\subsection{Does This Deny Gravitational Waves?}

No.

Gravitational waves are
measurable dynamical features
of spacetime geometry.

Their physical significance
is expressed through effects
such as changes in the relative
separation of freely falling
test masses.

The relational interpretation
does not deny those effects.

It rejects the additional
image of geometry as a material
container surrounding the
universe.

\subsection{Is This Merely a Change of Words?}

The distinction is conceptual,
but not empty.

The container interpretation
invites questions about
external location, enclosing
boundaries, and what lies
beyond the universe.

The relational interpretation
does not require those
additional structures.

It therefore changes the
ontological commitments
associated with the same
physical measurements.

An ontological clarification
need not alter established
equations to be meaningful.

It should, however, be
distinguished from a new
empirically confirmed theory.

\subsection{Could a Fundamental Space Simply Exist?}

A substantivalist may propose
a fundamental spatial or
spacetime structure that
does not require an enclosing
container.

That position avoids the
specific infinite-regress
argument developed here.

The regress argument therefore
does not logically eliminate
all forms of substantivalism.

The relational position
instead offers an economical
interpretation that does
not introduce independent
spatial substance beyond
the geometric relationships
required by physical theory.

\section{Discussion}

The ordinary language of
physics frequently treats
space as though it were
a physical location.

Objects are said to occupy
space, move through space,
and exist in space.

These expressions are useful
for calculations and everyday
communication.

Their usefulness does not
establish a literal container
ontology.

The distinction becomes
especially important when
ordinary spatial language
is applied to the universe
as a whole.

The universe does not
require an external address.

Its measurable geometry
is sufficient to define
relationships among its
constituents.

The conceptual error occurs
when the relational meaning
of location is transformed
into a demand that the
entire relational structure
must itself be located.

That demand produces
unnecessary explanatory
questions.

Where is the container?

What contains it?

What material forms its
boundary?

What lies outside?

These questions arise
from the analogy, not
from the measurements.

Removing the analogy
removes the regress.

The result is a simpler
ontological description
that preserves measurable
distance and intrinsic
geometry while dispensing
with an independently
existing spatial box.

\section{Conclusion}

Space is not a box
inside another box.

It is not a holodeck
sitting within a larger
physical room.

A material container
already presupposes the
spatial relationships
it is supposed to explain.

An infinite sequence
of enclosing containers
does not solve that
explanatory problem.

The relational alternative
begins with the measurable
relationships themselves.

Distance describes
separation.

Geometry describes
the organization of
separations and intervals.

Location identifies
relationships relative
to physical reference
systems.

None of these concepts
requires the universe
to occupy an external
physical container.

The conclusion is
therefore not that
distance is unreal,
nor that geometry
is dispensable.

It is that spatial
relationships need not
be reified into an
independent substance.

\begin{center}
\large
\textbf{
Space is not something
that exists somewhere.
}
\end{center}

\begin{center}
\large
\textbf{
Space is the relational
structure by which
``somewhere'' is defined.
}
\end{center}

\section*{Author's Note}

This paper develops an
ontological interpretation
of space rather than a
replacement for established
relativistic geometry.

Its central argument is
intended to stand independently
of the author's Timeless
Light Model.

The relational thesis
and its implications for
TLM are therefore
distinguished throughout.

\begin{thebibliography}{99}

\bibitem[Alexander(1956)]{alexander1956}
Alexander, H. G. (Ed.).
\emph{The Leibniz--Clarke Correspondence}.
Manchester University Press, 1956.

\bibitem[Barbour(1982)]{barbour1982}
Barbour, J. B.
Relational concepts of space and time.
\emph{British Journal for the Philosophy
of Science}, 33(3), 251--274, 1982.
\href{https://doi.org/10.1093/bjps/33.3.251}
{doi:10.1093/bjps/33.3.251}.

\bibitem[Earman(1989)]{earman1989}
Earman, J.
\emph{World Enough and Space-Time:
Absolute versus Relational Theories
of Space and Time}.
MIT Press, 1989.

\bibitem[Einstein(1916)]{einstein1916}
Einstein, A.
Die Grundlage der allgemeinen
Relativit\"atstheorie.
\emph{Annalen der Physik},
354(7), 769--822, 1916.
\href{https://doi.org/10.1002/andp.19163540702}
{doi:10.1002/andp.19163540702}.

\bibitem[Rovelli(2004)]{rovelli2004}
Rovelli, C.
\emph{Quantum Gravity}.
Cambridge University Press, 2004.

\bibitem[Wald(1984)]{wald1984}
Wald, R. M.
\emph{General Relativity}.
University of Chicago Press, 1984.

\end{thebibliography}

\end{document}
